
% XX大学博士毕业论文LaTeX模板v2.0.0
% 原作：白鸽坐飞机（知乎、简书和B站同号）。原格式可以认为是惨不忍堵 (～￣▽￣)～ 
% 改良：项目级重构版本
% 需要xelatex→bibtex→xelatex→xelatex进行完全编绎
% 本tex的内容可按需修改，但为了不打破依赖关系，本tex的布局不可轻易改动，特别是页脚部分！

%%%————————————————————————————————全局环境

\RequirePackage{xeCJK}
\RequirePackage{fontspec}%字体的包
\RequirePackage{bensz-fonts}
\newcommand{\tnewroman}{\BenszFontsUseTimesFallback}
\newcommand{\adobegothicstdb}{\fontspec{Adobe Gothic Std B}}
\newcommand{\Centaur}{\fontspec{Centaur}}
\BenszFontsSetMainTimesFallback %fontspec下这个命令设置全局默认字体
% \setmainfont{\songti\zihao{-4}} %fontspec下这个命令设置全局默认字体

%%%————————————————————————————————目录设置

% LaTeX目录设置教程：
% 用latex写毕业论文--用titletoc包重设目录格式： https://blog.csdn.net/G_Barble/article/details/111500095
% LaTeX目录格式控制：https://www.jianshu.com/p/d906e52b851d

\RequirePackage{titlesec,titletoc} %更改目录的包
\RequirePackage{afterpage}%更改目录计数

\renewcommand\contentsname{目\hspace{1em}录}

% \ctexset用于编辑section的样式，而非目录的样式。
\ctexset{
	section = {
		number = 第\chinese{section}章,
		format = \zihao{2}\bfseries\heiti\centering,
	},
	subsection = {
		number = \arabic{section}.\arabic{subsection},
		format = \zihao{-2}\bfseries\heiti,
		beforeskip = 0.18em,
		afterskip = 0.12em,
	},
	subsubsection = {
		number = \arabic{section}.\arabic{subsection}.\arabic{subsubsection},
		format = \zihao{-3}\bfseries\heiti,
	},
	paragraph = {
		number = \arabic{section}.\arabic{subsection}.\arabic{subsubsection}.\arabic{paragraph},
		format = \zihao{4}\bfseries\heiti,
	}
}

% 直接用 titlesec 精确压紧小节标题上下留白；仅靠 ctexset 的
% beforeskip/afterskip 在当前模板里体感不够明显。
\titlespacing*{\subsection}{0pt}{4pt}{2pt}

\titlecontents{section}[4em]
{\bfseries\songti\zihao{-4}\vspace{0.4em}}
{\contentslabel{3.8em}}
{\hspace*{-3.8em}}
{~\titlerule*[0.6pc]{$.$}~\contentspage}

\titlecontents{subsection}[3.5em]
{\songti\zihao{-4}\vspace{0.2em}}
{\contentslabel{2em}}
{\hspace*{-2em}}
{~\titlerule*[0.6pc]{$.$}~\contentspage} 

\titlecontents{subsubsection}[5.5em]
{\songti \zihao{-4}\vspace{0.2em}}
{\contentslabel{2.5em}}
{\hspace*{-2em}}
{~\titlerule*[0.8pc]{$.$}~\contentspage}

\setcounter{tocdepth}{3}% 设置目录级数
\setcounter{secnumdepth}{4}% 设置在几级目录前标记序号

% \titlecontents用于编辑目录的样式，而非section的样式。
% \titlecontents{section}[4em]%标题位置左间距
% {\bfseries\songti\zihao{-4}\vspace{10pt}}%设置标题字体是宋体；字号小四；与上一个标题垂直距离
% {\contentslabel{3.8em}}%设置标题标志的格式，如序号格式、序号宽度、序号与标题内容之间的间距
% {\hspace*{-2em}}%设置无序号标题的格式，如字体、字体尺寸、位置等。该参数可以空置
% {~\titlerule*[0.6pc]{$.$}~\contentspage}%设置标题与页码之间的指引线样式以及页码的格式，该参数如果空置，标题将无指引线


%%%————————————————————————————————图片与表格

% 对于表格，推荐使用Excel的插件Excel2LaTeX自动生成LaTeX代码。详见： https://www.zhihu.com/question/307970489/answer/2305355098

\RequirePackage{graphicx} %插图宏包
\RequirePackage{float} %设置图片浮动位置的宏包
\RequirePackage{subfigure} %插入多图时用子图显示的宏包
\RequirePackage[section]{placeins} %避免浮动体跨过\section
\RequirePackage{enumerate} %设置列表环境的包

\RequirePackage{caption}
\captionsetup[figure]{name={图},font={footnotesize,stretch=1.25},labelsep=period,labelfont=bf, singlelinecheck=off, justification=centering} 
\captionsetup[table]{name={表},font={footnotesize,stretch=1.25},labelsep=period,labelfont=bf, singlelinecheck=off, justification=centering} 
% \RequirePackage[font=small,labelfont=bf,labelsep=none]{caption}
% \captionsetup[table]{
%   %labelsep=newline,%换行
%   singlelinecheck=false,%居左
% }

% 将每页的图片置顶展示（而不是居中）
\makeatletter	
	\setlength{\@fptop}{0pt}
\makeatother

\RequirePackage{booktabs} %表格
\RequirePackage{tabularx} %表格宽度
\RequirePackage{multirow} %多行合并

% [latex Longtable使用指南-字体调节，与table不同\_易烊千蝈的博客-CSDN博客](https://blog.csdn.net/weixin_39490300/article/details/128184293)
% [LaTeX使用笔记：长表格longtable（附实例） – Spark & Shine](http://sparkandshine.net/latex-use-notes-longtable-with-examples/)
\RequirePackage{longtable} % latex 多页显示同一表格
\RequirePackage{tabu} % 使用longtabu（tabularx + longtable）将长文本在单元格多行显示。

%%%————————————————————————————————段落与行

% LaTeX 段落：段落缩进、段落间距、行距：https://blog.csdn.net/xovee/article/details/119654826
\RequirePackage{setspace}%设置行距的包
%\setlength{\baselineskip}{1.6em}%2倍行距
\linespread{1.3} % 设置行距为1.5倍行距。1.0，单倍行距；1.3，一点五倍行距；1.6，双倍行距。
\setlength{\parskip}{0.5em} % 将段落之间的间距设置为0.5em

% \RequirePackage{indentfirst}%首行缩进的包
% \setlength{\parindent}{2em}%首行缩进2。不过我一般比较喜欢局部而非全局定义。

\RequirePackage{geometry}%设置页边距的包
\geometry{left=2.5cm,right=2.5cm,top=3.5cm,bottom=2.5cm} % 设置页边距
\raggedbottom % 避免为齐底强行拉伸标题前后的竖向留白

% 参考：https://blog.csdn.net/jzwong/article/details/99964922
\RequirePackage{ragged2e} % 正文两端对齐所需的包。在需要两端对齐的文字前面添加\justifying即可。

\RequirePackage{enumitem} % 增强自定义Bullet（即各种List前的那个标记）


%%%————————————————————————————————代码与公式

% 编辑数学公式推荐使用“在线LaTeX公式编辑器”： https://www.latexlive.com/

\RequirePackage{bm} %加粗公式。使用方法为：\bm{$abc$}
\RequirePackage{listings} %代码块
\RequirePackage{xcolor} %代码块
\renewcommand{\lstlistingname}{Code}
% \lstset{
% 	% language=bash,  %代码语言使用的是matlab
% 	frame=shadowbox, %把代码用带有阴影的框圈起来
% 	rulesepcolor=\color{red!20!green!20!blue!20},%代码块边框为淡青色
% 	keywordstyle=\color{blue!90}\bfseries, %代码关键字的颜色为蓝色，粗体
% 	commentstyle=\color{red!10!green!70}\textit,% 设置代码注释的颜色
% 	showstringspaces=true,%不显示代码字符串中间的空格标记
% 	numbers=left, % 显示行号
% 	numberstyle=\tiny,    % 行号字体
% 	stringstyle=\ttfamily, % 代码字符串的特殊格式
% 	breaklines=false, %对过长的代码自动换行
% 	extendedchars=false,  %解决代码跨页时，章节标题，页眉等汉字不显示的问题
% 	escapebegin=\begin{CJK*},escapeend=\end{CJK*},% 代码中出现中文必须加上，否则报错
% 	texcl=true
% }
% 教程见：
% https://www.youtube.com/watch?v=qxkQgG1Y0bY
% https://zhuanlan.zhihu.com/p/65441079
\lstdefinestyle{codestyle01}{
	backgroundcolor=\color{gray!12},
	%basicstyle=\ttfamily\small,
	basicstyle=\zihao{-5}\ttfamily,
	commentstyle=\color{green!60!black},
	keywordstyle=\color{magenta},
	stringstyle=\color{blue!50!red},
	showstringspaces=false,
	% captionpos=
	numbers=left,
	numberstyle=\footnotesize\color{gray},
	numbersep=10pt,
	stepnumber=2, %左边数字的步幅
	tabsize=4,
	frame=TB %顶部和底部有两条竖线。
	% frame=tblr
	% frame=L,
	% framerule=1pt,
	% rulecolor=\color{red}
}

%%%————————————————————————————————引用

\RequirePackage{cite}%设置参考文献的包
%参考https://blog.csdn.net/wy_wy12/article/details/95604258
%plain，按字母的顺序排列，比较次序为作者、年度和标题.
%unsrt，样式同plain，只是按照引用的先后排序.
%alpha，用作者名首字母+年份后两位作标号，以字母顺序排序.
%abbrv，类似plain，将月份全拼改为缩写，更显紧凑.
%ieeetr，国际电气电子工程师协会期刊样式.
%acm，美国计算机学会期刊样式.
%siam，美国工业和应用数学学会期刊样式.
%apalike，美国心理学学会期刊样式.
% 参考文献misc无法正常引用
% https://gitlab.com/sysu-gitlab/latex-group/thesis/-/issues/96
\makeatletter
	\def\@cite#1#2{\textsuperscript{[{#1\if@tempswa , #2\fi}]}} % 参考文献标注置于右上角
\makeatother


% \RequirePackage{hyperref}%内部交叉引用
% \hypersetup{hidelinks}%去除参考文献的绿框
\RequirePackage[
	colorlinks,
	urlcolor=blue, % 超链接的颜色。默认是芭比粉
	linkcolor=black,
	anchorcolor=black,
	citecolor=black,
	CJKbookmarks=True
]{hyperref}


%%%————————————————————————————————页脚页眉

%页眉参考http://www.ctex.org/documents/packages/layout/fancyhdr.htm
\RequirePackage{fancyhdr}  %设置页眉页码的包
% 字体大小 \tiny \scriptsize \footnotesize \small \normalsize \large \Large \LARGE \huge \Huge 或者 \zihao{number}{加入文字}
\pagestyle{fancy}
\lhead{\footnotesize 中山大学博士学位论文} %页眉左边 
\chead{ } %页眉中间
\rhead{\footnotesize \titleZh} %页眉右边显示论文中文标题
%\lfoot{From: K. Grant} %页脚右边是From: K. Grant
\cfoot{\thepage} %页脚中间是本页页码
%\rfoot{To XXX} %页脚右边是To XXX
\renewcommand{\headrulewidth}{0.4pt} %页眉划线
%\renewcommand{\footrulewidth}{0.4pt}%页脚划线
\RequirePackage{pgf} % 密级右上角

%%%————————————————————————————————其它

%-----解决\headheight is too small的warning
\setlength{\headheight}{15pt}

%-----输出docx
% PowerShell: 
% pandoc --version % pandoc.exe 2.17.1.1
% cd E:\iProjects\Manuscripts\thesis_sysu; 
% pandoc SYSU.palte.tex  --filter pandoc-crossref --citeproc --csl springer-basic-note.csl  --bibliography=Graduate.bib -M reference-section-title=Reference -o SYSU.palte_pandoc.docx

% (.*) 删除备注？

%------LaTeX入门和学习教程
% https://nextcloud.hwb0307.com/s/t36QzwSnNCEBoR2。来自LaTeX学习交流群（qq群1019471316）。

%------模块化设计思想
%  Splitting a Large Document into Several Files: https://www.dickimaw-books.com/latex/thesis/html/include.html
